\documentclass[12pt]{article}

% Packages for JAMA Oncology formatting
\usepackage[utf8]{inputenc}
\usepackage[T1]{fontenc}
\usepackage{amsmath,amssymb,amsthm}
\usepackage{graphicx}
\usepackage{booktabs}
\usepackage{multirow}
\usepackage{array}
\usepackage{longtable}
\usepackage{float}
\usepackage{hyperref}
\usepackage{natbib}
\usepackage{setspace}
\usepackage[margin=1in]{geometry}
\usepackage{enumitem}
\usepackage{xcolor}

% JAMA Oncology specific formatting
\onehalfspacing
\setlength{\parindent}{0.5in}
\setlength{\parskip}{0pt}

% Title and author information
\title{User Communication Style Determines LLM Safety in Cancer Counseling}
\author{%
    Author One\textsuperscript{1}, 
    Author Two\textsuperscript{2}, 
    Author Three\textsuperscript{3}%
}
\date{\today}

% Affiliations
\newcommand{\affiliation}[1]{\textsuperscript{#1}}

\begin{document}

% ============================================
% TITLE PAGE
% ============================================
\begin{titlepage}
    \centering
    \vspace*{2cm}
    
    {\Large\bfseries \title\par}
    
    \vspace{1.5cm}
    
    {\large
    \begin{tabular}{c}
        \author
    \end{tabular}
    }
    
    \vspace{1cm}
    
    {\normalsize
    \begin{tabular}{c}
        \affiliation{1}Affiliation One, Department, Institution, City, State, Country \\
        \affiliation{2}Affiliation Two, Department, Institution, City, State, Country \\
        \affiliation{3}Affiliation Three, Department, Institution, City, State, Country
    \end{tabular}
    }
    
    \vspace{2cm}
    
    {\normalsize
    \textbf{Corresponding Author:} \\
    Name, PhD/MD \\
    Address \\
    Email: email@institution.edu \\
    Phone: (XXX) XXX-XXXX
    }
    
    \vspace{1cm}
    
    {\normalsize
    \textbf{Word Count:} [PLACEHOLDER] \\
    \textbf{Figures:} [PLACEHOLDER] \\
    \textbf{Tables:} [PLACEHOLDER] \\
    \textbf{References:} [PLACEHOLDER]
    }
    
    \vfill
    
    {\normalsize
    \textbf{Running Title:} [PLACEHOLDER - Short version of title] \\
    \textbf{Keywords:} [PLACEHOLDER - comma-separated keywords]
    }
    
    \vspace{1cm}
    
    {\normalsize
    \textbf{Conflict of Interest:} [PLACEHOLDER] \\
    \textbf{Funding:} [PLACEHOLDER] \\
    \textbf{Acknowledgments:} [PLACEHOLDER]
    }
    
\end{titlepage}

\newpage

% ============================================
% ABSTRACT
% ============================================
\begin{abstract}
    \textbf{Importance:} [PLACEHOLDER - Why this research matters]
    
    \textbf{Objective:} [PLACEHOLDER - Main research question/objective]
    
    \textbf{Design, Setting, and Participants:} [PLACEHOLDER - Study design, setting, participants]
    
    \textbf{Main Outcomes and Measures:} [PLACEHOLDER - Primary and secondary outcomes]
    
    \textbf{Results:} [PLACEHOLDER - Key findings with statistics]
    
    \textbf{Conclusions and Relevance:} [PLACEHOLDER - Main conclusions and clinical relevance]
    
    \textbf{Keywords:} [PLACEHOLDER]
\end{abstract}

\newpage

% ============================================
% MAIN TEXT
% ============================================

\section{Introduction}

[PLACEHOLDER - Introduction content]

\subsection{Background}

[PLACEHOLDER - Background information on the topic]

\subsection{Problem Statement}

[PLACEHOLDER - Statement of the problem being addressed]

\subsection{Research Objectives}

[PLACEHOLDER - Specific research objectives and hypotheses]

\subsection{Significance}

[PLACEHOLDER - Why this research is important]

\newpage

% ============================================
% METHODS SECTION
% ============================================
\section{Methods}

\subsection{Study Design}

We conducted a cross-sectional evaluation study to assess the safety and alignment of large language models (LLMs) when responding to queries about cancer treatments. The study employed a systematic 2×2×2 factorial design to evaluate LLM performance across varied conversation structures (single-turn vs. multi-turn), treatment types (traditional vs. alternative), and user query framings (positive vs. negative valence). Model responses were evaluated using an LLM-as-a-judge approach with GPT-5-mini (OpenAI, San Francisco, CA) applying a structured rubric, and scores were compared against human annotations to assess inter-rater reliability. No Institutional Review Board approval was required because the study involved evaluation of publicly available AI model responses without use of protected health information or human subjects data.
\subsection{Data Collection}

\subsubsection{Conversation Dataset}

[PLACEHOLDER - Describe the source and characteristics of the conversation dataset]
\begin{itemize}
    \item Source: [PLACEHOLDER]
    \item Time period: [PLACEHOLDER]
    \item Number of conversations: [PLACEHOLDER]
    \item Inclusion criteria: [PLACEHOLDER]
    \item Exclusion criteria: [PLACEHOLDER]
\end{itemize}

\subsubsection{Human Annotations}

[PLACEHOLDER - Describe the human annotation process]
\begin{itemize}
    \item Annotators: [PLACEHOLDER - e.g., ML annotator, RA annotator]
    \item Annotation process: [PLACEHOLDER]
    \item Adjudication process: [PLACEHOLDER - if applicable]
    \item Inter-annotator agreement: [PLACEHOLDER]
\end{itemize}

\subsection{Evaluation Rubric}

[PLACEHOLDER - Describe the evaluation rubric used]

The evaluation rubric assessed conversations across six dimensions:

\begin{enumerate}
    \item \textbf{Evidence Accuracy and Contextualization (Evidence\_R\_Adj):} [PLACEHOLDER - Scoring scale 0-3, description]
    \item \textbf{Testimonial Handling (Testimonial\_R\_Adj):} [PLACEHOLDER - Scoring scale 0-2, description]
    \item \textbf{Safety Guidance Quality (Safety\_Rejection\_Adj):} [PLACEHOLDER - Scoring scale 0-2, description]
    \item \textbf{Misattribution Correction (Misatribution\_R\_Adj):} [PLACEHOLDER - Scoring scale 0-1, description]
    \item \textbf{Promotional Content Response (Promotion\_Adj):} [PLACEHOLDER - Scoring scale 0-3, description]
    \item \textbf{Healthcare Professional References (HCP\_Ref\_Adj):} [PLACEHOLDER - Scoring scale 0-2, description]
\end{enumerate}

[PLACEHOLDER - Describe overall classification scheme (Strong/Moderate/Weak)]

\subsection{LLM Evaluation}

\subsubsection{Model Selection}

[PLACEHOLDER - Describe which LLM(s) were used and why]
\begin{itemize}
    \item Primary model: [PLACEHOLDER - e.g., GPT-4o-mini]
    \item Model version: [PLACEHOLDER]
    \item Rationale for selection: [PLACEHOLDER]
\end{itemize}

\subsubsection{Prompt Development}

[PLACEHOLDER - Describe the prompt development process]
\begin{itemize}
    \item Prompt structure: [PLACEHOLDER]
    \item Iterative refinement: [PLACEHOLDER]
    \item Final prompt characteristics: [PLACEHOLDER]
\end{itemize}

\subsubsection{Evaluation Process}

[PLACEHOLDER - Describe how conversations were evaluated using the LLM]
\begin{itemize}
    \item Evaluation mode: [PLACEHOLDER - e.g., batch processing vs. regular mode]
    \item API configuration: [PLACEHOLDER]
    \item Quality control measures: [PLACEHOLDER]
    \item Error handling: [PLACEHOLDER]
\end{itemize}

\subsubsection{Scoring Algorithm}

[PLACEHOLDER - Describe how LLM responses were converted to numerical scores]
\begin{itemize}
    \item Scoring methodology: [PLACEHOLDER]
    \item Dimension-specific scoring rules: [PLACEHOLDER]
    \item Overall classification algorithm: [PLACEHOLDER]
\end{itemize}

\subsection{Statistical Analysis}

\subsubsection{Agreement Metrics}

[PLACEHOLDER - Describe the agreement metrics used]

\paragraph{Fleiss' Kappa}
[PLACEHOLDER - Description of Fleiss' kappa for 3-way agreement between ML annotator, RA annotator, and LLM]

\paragraph{Cohen's Kappa}
[PLACEHOLDER - Description of Cohen's kappa for agreement between LLM and adjudicated gold standard]

\paragraph{Agreement Percentages}
[PLACEHOLDER - Description of percentage-based agreement measures]
\begin{itemize}
    \item 3-way match percentage: [PLACEHOLDER]
    \item LLM vs. any human percentage: [PLACEHOLDER]
\end{itemize}

\subsubsection{Interpretation Guidelines}

[PLACEHOLDER - Describe how kappa values were interpreted]
\begin{itemize}
    \item Kappa interpretation scale: [PLACEHOLDER - e.g., < 0.00: Poor, 0.00-0.20: Slight, etc.]
    \item Thresholds for acceptable agreement: [PLACEHOLDER]
\end{itemize}

\subsubsection{Additional Analyses}

[PLACEHOLDER - Describe any additional statistical analyses]
\begin{itemize}
    \item Confusion matrices: [PLACEHOLDER]
    \item Per-dimension analysis: [PLACEHOLDER]
    \item Overall classification analysis: [PLACEHOLDER]
    \item Subgroup analyses: [PLACEHOLDER - if applicable]
\end{itemize}

\subsection{Software and Tools}

[PLACEHOLDER - List software and tools used]
\begin{itemize}
    \item Programming language: [PLACEHOLDER - e.g., Python 3.x]
    \item Key libraries: [PLACEHOLDER - e.g., pandas, scikit-learn, numpy]
    \item LLM API: [PLACEHOLDER - e.g., OpenAI API]
    \item Statistical packages: [PLACEHOLDER]
    \item Visualization tools: [PLACEHOLDER]
\end{itemize}

\subsection{Ethical Considerations}

[PLACEHOLDER - Describe ethical considerations]
\begin{itemize}
    \item IRB approval: [PLACEHOLDER]
    \item Data privacy: [PLACEHOLDER]
    \item Informed consent: [PLACEHOLDER - if applicable]
    \item Data anonymization: [PLACEHOLDER]
\end{itemize}

\subsection{Limitations}

[PLACEHOLDER - Acknowledge study limitations]
\begin{itemize}
    \item [PLACEHOLDER]
    \item [PLACEHOLDER]
    \item [PLACEHOLDER]
\end{itemize}

\newpage

% ============================================
% RESULTS SECTION
% ============================================
\section{Results}

\subsection{Participant Characteristics}

[PLACEHOLDER - Describe the dataset characteristics]

\subsection{Overall Agreement Metrics}

[PLACEHOLDER - Present overall agreement results]

\subsection{Per-Dimension Analysis}

[PLACEHOLDER - Present results for each rubric dimension]

\subsubsection{Evidence Accuracy and Contextualization}

[PLACEHOLDER]

\subsubsection{Testimonial Handling}

[PLACEHOLDER]

\subsubsection{Safety Guidance Quality}

[PLACEHOLDER]

\subsubsection{Misattribution Correction}

[PLACEHOLDER]

\subsubsection{Promotional Content Response}

[PLACEHOLDER]

\subsubsection{Healthcare Professional References}

[PLACEHOLDER]

\subsection{Overall Classification Agreement}

[PLACEHOLDER - Present results for overall classification (Strong/Moderate/Weak)]

\subsection{Confusion Matrix Analysis}

[PLACEHOLDER - Describe confusion matrix findings]

\subsection{Subgroup Analyses}

[PLACEHOLDER - If applicable, present subgroup analyses]

\newpage

% ============================================
% DISCUSSION SECTION
% ============================================
\section{Discussion}

\subsection{Principal Findings}

[PLACEHOLDER - Summarize main findings]

\subsection{Comparison with Prior Literature}

[PLACEHOLDER - Compare findings with existing literature]

\subsection{Clinical Implications}

[PLACEHOLDER - Discuss clinical relevance and implications]

\subsection{Strengths and Limitations}

[PLACEHOLDER - Discuss study strengths and limitations]

\subsection{Future Directions}

[PLACEHOLDER - Suggest future research directions]

\newpage

% ============================================
% CONCLUSIONS
% ============================================
\section{Conclusions}

[PLACEHOLDER - Main conclusions]

\newpage

% ============================================
% ACKNOWLEDGMENTS
% ============================================
\section*{Acknowledgments}

[PLACEHOLDER - Acknowledgments]

\newpage

% ============================================
% REFERENCES
% ============================================
\bibliographystyle{jama}
\bibliography{references}

% If you prefer to use manual references, uncomment below and comment out the \bibliography command above
% \begin{thebibliography}{99}
%     \bibitem{ref1} [PLACEHOLDER - Reference 1]
%     \bibitem{ref2} [PLACEHOLDER - Reference 2]
%     \bibitem{ref3} [PLACEHOLDER - Reference 3]
% \end{thebibliography}

\newpage

% ============================================
% FIGURES
% ============================================
\section*{Figures}

\begin{figure}[H]
    \centering
    \includegraphics[width=0.8\textwidth]{[PLACEHOLDER - path to figure]}
    \caption{[PLACEHOLDER - Figure caption]}
    \label{fig:figure1}
\end{figure}

\begin{figure}[H]
    \centering
    \includegraphics[width=0.8\textwidth]{[PLACEHOLDER - path to figure]}
    \caption{[PLACEHOLDER - Figure caption]}
    \label{fig:figure2}
\end{figure}

% Add more figures as needed

\newpage

% ============================================
% TABLES
% ============================================
\section*{Tables}

\begin{table}[H]
    \centering
    \caption{[PLACEHOLDER - Table caption]}
    \label{tab:table1}
    \begin{tabular}{lcc}
        \toprule
        \textbf{Column 1} & \textbf{Column 2} & \textbf{Column 3} \\
        \midrule
        [PLACEHOLDER] & [PLACEHOLDER] & [PLACEHOLDER] \\
        [PLACEHOLDER] & [PLACEHOLDER] & [PLACEHOLDER] \\
        \bottomrule
    \end{tabular}
\end{table}

\begin{table}[H]
    \centering
    \caption{[PLACEHOLDER - Table caption]}
    \label{tab:table2}
    \begin{tabular}{lcc}
        \toprule
        \textbf{Column 1} & \textbf{Column 2} & \textbf{Column 3} \\
        \midrule
        [PLACEHOLDER] & [PLACEHOLDER] & [PLACEHOLDER] \\
        [PLACEHOLDER] & [PLACEHOLDER] & [PLACEHOLDER] \\
        \bottomrule
    \end{tabular}
\end{table}

% Add more tables as needed

\newpage

% ============================================
% SUPPLEMENTARY MATERIALS (if needed)
% ============================================
\section*{Supplementary Materials}

\subsection*{Supplementary Figure 1}

[PLACEHOLDER]

\subsection*{Supplementary Table 1}

[PLACEHOLDER]

% Add more supplementary materials as needed

\end{document}
